\section{Conclusion}\label{sec:discussion}
Compact support allows the singular solution to be placed inside a small
region, and parabolic scaling makes its force small below the critical
orders. A cutoff of a local vector potential removes the reference velocity
in that region while preserving incompressibility. The resulting exact
construction proves density for every fixed admissible smooth initial
velocity in the force topologies specified above.

For zero initial velocity, the converse follows from regularity estimates.
On both $\TT$ and $\R^3$, sufficiently small $L^1_tH^{1/2}_x$ force gives
global regularity. On $\R^3$, sufficiently small $L^2_tH^{-1/2}_x$ force
gives regularity on each prescribed finite interval; the permitted radius
depends on that interval. These estimates give the sharp thresholds in
Theorems~\ref{thm:main} and~\ref{thm:Rmain}. The whole-space proof controls
low frequencies separately, both when scaling negative Sobolev norms and
when obtaining the $H^2$ integrability needed for continuation.

The approximation preserves more than the initial velocity. Near every
regular reference, it preserves the earlier history and converges strongly
in energy and dissipation. On $\R^3$ it also preserves all velocity and
force cell averages for every $0\le t<T$ on any fixed finite family of
Cartesian grids, when the insertion is placed inside a common cell.
The periodic and bounded-domain constructions give prescribed finite
collections of singular regions. The whole-space compact solution also
admits the infinite-dimensional affine family of velocities in
Proposition~\ref{prop:affine}, with corresponding changes in the force.

The quantifiers and topologies determine the scope of these conclusions.
The force may depend on the initial velocity; no classification of singular
initial velocities for a single prescribed force follows. Exact blowup at
$T$ is obtained near references regular through $T$, whereas the general
density statement concerns breakdown by $T$. Small integrated force norms
coexist with increasing force amplitudes and derivative norms. Thus density
alone gives neither a probability of blowup nor stability under arbitrary
perturbations, and it does not imply density within a fixed spectral
truncation or a prescribed finite-dimensional force family.

Finally, the energy estimates concern the approach to the singular time.
They do not select a continuation afterward or assert nonexistence or
nonuniqueness of Leray--Hopf continuations. The grid statement concerns
specified observations; the construction is not a computational fluid
dynamics experiment.
